\documentclass[10pt,a4paper]{amsart}
\usepackage[margin=19mm]{geometry}
\usepackage{amsmath,amssymb,mathtools}
\usepackage[scr=rsfs,cal=euler]{mathalfa}
\usepackage{xfrac}
\usepackage[unicode,hidelinks,bookmarks=false]{hyperref}
\newcommand{\ie}{i.e.\ }
\newcommand{\cf}{cf.\ }
\newcommand{\resp}{respectively\ }
\newcommand{\Subsection}{\subsection}
\providecommand{\nobreakdash}{\nobreak\hbox{-}}
\setlength{\emergencystretch}{2em}
\multlinegap0pt
\usepackage{bm}
\usepackage{bbm}
\usepackage{dsfont}
\usepackage{xspace}
\usepackage{cancel}
\usepackage{mathtools}
\usepackage{amsthm}
\makeatletter
\@ifundefined{proposition}{\newtheorem{proposition}{Proposition}}{}
\makeatother
\newcommand{\eps}{\varepsilon}
\title[Balanced Dynamics in Strongly Coupled Networks]{Balanced Dynamics in Strongly Coupled Networks}
\author{Cristobal Quininao, Jonathan Touboul}
\date{Source: arXiv:2501.11769v2 (CC BY 4.0)}
\begin{document}
\maketitle

\noindent\emph{Source and licence.} Balanced Dynamics in Strongly Coupled Networks. Version arXiv:2501.11769v2 (CC BY 4.0), distributed under the Creative Commons Attribution 4.0 International licence, as recorded in the arXiv licence metadata for this exact version and shown on the abstract page https://arxiv.org/abs/2501.11769v2. Full rights notice: licenses/arxiv-2501.11769-CC-BY-4.0.txt.\par\smallskip

\noindent\textit{Editorial scope.} Counted unit: the Proposition whose source label is \texttt{prop:regularity\_space} in the article (source0.tex lines 1460--1466, source label \texttt{prop:regularity\_space}), delivered together with the article's own complete original proof of it (source0.tex lines 1468--1567, one \texttt{proof} environment of 100 lines). No mathematics is written, repaired or re-derived: statement and proof are the article's own text line for line. Required context retained uncounted: prop:supersol (lines 1384--1395). Every definition, display and label of those lines is retained unchanged. Omitted from this package and not redistributed: 1--1383, 1396--1459, 1467--1467, 1568--2338. Nothing inside a retained excerpt is omitted or altered: every retained body is delivered exactly as the article's lines are written, so no omission receipt of any kind accompanies this package. Citations: the retained excerpts carry no citation command at all, so no bibliography entry of the article is transcribed and no reference is redistributed. Reference adaptation: none was needed and none was made; every reference inside the delivered unit resolves inside it, so no unresolved reference or citation is produced. Printed slips kept literal: no printed word, symbol, hypothesis or number of the counted statement or of its proof was altered. Layout and class: the article's own class is \texttt{siamart190516} with packages including arydshln, cite, lipsum, amsfonts, graphicx, epstopdf, mathtools, bbm, enumerate, subfig, algorithmic, cancel, comment, hyperref; the wrapper is the lane's pinned \texttt{amsart} editorial wrapper at 10pt on A4, which restates the theorem-like environments the retained text needs and adds the article's own macro definitions that the retained text needs (\texttt{\textbackslash eps}), with extra wrapper packages (bm, bbm, dsfont, xspace, cancel, mathtools). The delivered numbering is the unit's own. Assets: no retained span contains a figure environment, a graphics-inclusion command, a PDF-page inclusion, a table or a TikZ picture; no image file is copied into this package, the package declares no asset dependency and no image file is redistributed with it. Licence: version arXiv:2501.11769v2 of this article is distributed under the Creative Commons Attribution 4.0 International licence, as recorded in the arXiv licence metadata for this exact version and shown on the abstract page https://arxiv.org/abs/2501.11769v2. A full rights notice is delivered at licenses/arxiv-2501.11769-CC-BY-4.0.txt. Characters: every retained excerpt is pure ASCII, so the delivered engine, XeLaTeX with its default Latin Modern text font, reports no missing character.
\begin{proposition}
\label{prop:supersol}
Assume that initial conditions $\phi_\eps(0,x)=\eps\log\mu_{\eps,0}(x)$ satisfy the conditions
$$
\sup_{0<\eps<1}\eps\log \mu_{0,\eps}(x) \leq -Ax^2+B,
$$
for some constants $A,B$ with $A$ positive. Then there is another set of positive constants $A',B',D'$ and $E'$ such that
$$
\phi_\eps(t,x)\leq -A'I_\eps(t)\Lambda(x)-\frac{\eps}2 B'x^2+D't+E',
$$ 
at all positive times.
\end{proposition}

\begin{proposition}
\label{prop:regularity_space}
Assume that $\eps\leq 1$ and let $F=\sqrt{E'+D'T}$ for $E'$ and $D'$ the same positive constants from Proposition~\ref{prop:supersol}, and let $T>0$ be an arbitrary time horizon. Then the map $w_\eps(t,x)=\sqrt{2F^2-\phi_\eps(t,x)}$ is well defined and there exists a constant $\theta(T)$independent of $\eps$ such that
$$
\left|\partial_x w_\eps(t,x)\right|\leq  \sqrt{\frac{1}{t\sigma^2}}+\theta(T).
$$
\end{proposition}

\begin{proof}
First, notice that proposition~\ref{prop:supersol} implies in particular that
$$
\phi_\eps(t,x)\leq D'T+E'\quad\Rightarrow\quad 2F^2-\phi_\eps(t,x)\geq D'T+E'>0,
$$
and thus ensures that $w_\eps$ is well-defined. Our proof of an upper bound relies on a well-chosen change of variable. In detail, let $\gamma$ be a smooth invertible map, with nonvanishing derivative, and define $\phi_\eps(t,x)=\gamma(w_\eps(t,x))$. Elementary calculus yields:
$$
\partial_t\phi_\eps = \gamma'(w_\eps)\partial_tw_\eps,\qquad \partial_x\phi_\eps = \gamma'(w_\eps)\partial_xw_\eps,
$$
and
$$
\partial^2_{xx}\phi_\eps=\gamma''(w_\eps)|\partial_xw_\eps|^2+\gamma'(w_\eps)\partial^2_{xx}w_\eps,
$$
implying that $w_\eps$ is a solution to
\begin{multline}
\partial_tw_\eps = -\frac{\left(\eps f'(x) - \alpha'(x)I_{\mu_\eps}(t) \right)}{\gamma'(w_\eps)}-  \left( f(x)- \eps^{-1}\alpha(x) I_{\mu_\eps}(t)\right)\partial_xw_\eps
\\
 + \frac{\sigma^2}{2}\left(\frac{\gamma'(w_\eps)}{\eps}+\frac{\gamma''(w_\eps)}{\gamma'(w_\eps)}\right) \vert \partial_x w_\eps \vert^2  + \frac {\sigma^2}{2} \partial_{xx}^2 w_\eps.
\label{eq:wepsilon}
\end{multline}
Define $p_\eps(t,x):=\partial_xw_\eps$, and take an extra spatial partial derivative on~\eqref{eq:wepsilon} to get
\begin{multline*}
\partial_tp_\eps- \frac {\sigma^2}{2} \partial_{xx}^2 p_\eps = -\frac{\left(\eps f''(x) - \alpha''(x)I_{\mu_\eps}(t) \right)}{\gamma'(w_\eps)}+\frac{\left(\eps f'(x) - \alpha'(x)I_{\mu_\eps}(t) \right)}{|\gamma'(w_\eps)|^2}\gamma''(w_\eps)\,p_\eps
\\
-  \left( f'(x)- \eps^{-1}\alpha'(x) I_{\mu_\eps}(t)\right)p_\eps-  \left( f(x)- \eps^{-1}\alpha(x) I_{\mu_\eps}(t)\right)\partial_xp_\eps
\\
 + \frac{\sigma^2}{2}\left(\frac{\gamma''(w_\eps)}{\eps}+\frac{\gamma'''(w_\eps)}{\gamma'(w_\eps)}-\frac{\gamma''(w_\eps)^2}{\gamma'(w_\eps)^2}\right)  p_\eps^3
 \\ 
 + \sigma^2\left(\frac{\gamma'(w_\eps)}{\eps}+\frac{\gamma''(w_\eps)}{\gamma'(w_\eps)}\right) p_\eps \cdot\partial_x p_\eps,
\end{multline*}
and multiplying by $p_\eps$ the equation becomes
\begin{multline*}
\partial_t\frac{|p_\eps|^2}{2}-\frac{\sigma^2}{2}|\partial_x|p_\eps||^2-\frac{\sigma^2}{2}|p_\eps|\partial^2_{xx}|p_\eps|+\frac{\sigma^2}{2}|\partial_xp_\eps|^2
=
\\
-\frac{\left(\eps f''(x) - \alpha''(x)I_{\mu_\eps}(t) \right)}{\gamma'(w_\eps)}\, p_\eps+\frac{\left(\eps f'(x) - \alpha'(x)I_{\mu_\eps}(t) \right)}{|\gamma'(w_\eps)|^2}\gamma''(w_\eps)\,|p_\eps|^2
\\
-  \left( f'(x)- \eps^{-1}\alpha'(x) I_{\mu_\eps}(t)\right)\,|p_\eps|^2-  \frac12\left( f(x)- \eps^{-1}\alpha(x) I_{\mu_\eps}(t)\right)\partial_x|p_\eps|^2
\\
 + \frac{\sigma^2}{2}\left(\frac{\gamma''(w_\eps)}{\eps}+\frac{\gamma'''(w_\eps)}{\gamma'(w_\eps)}-\frac{\gamma''(w_\eps)^2}{\gamma'(w_\eps)^2}\right) |p_\eps|^4
 \\ 
 + \frac{\sigma^2}{2}\left(\frac{\gamma'(w_\eps)}{\eps}+\frac{\gamma''(w_\eps)}{\gamma'(w_\eps)}\right) p_\eps \cdot\partial_x |p_\eps|^2.
\end{multline*}
We find that:
\begin{multline*}
\partial_t|p_\eps|-\frac{\sigma^2}{2}\partial^2_{xx}|p_\eps|
\leq 
\\
-\frac{\left(\eps f''(x) - \alpha''(x)I_{\mu_\eps}(t) \right)}{\gamma'(w_\eps)}\, \frac{p_\eps}{|p_\eps|}+\frac{\left(\eps f'(x) - \alpha'(x)I_{\mu_\eps}(t) \right)}{|\gamma'(w_\eps)|^2}\gamma''(w_\eps)\,|p_\eps|
\\
-  \left( f'(x)- \eps^{-1}\alpha'(x) I_{\mu_\eps}(t)\right)|p_\eps|-  \left( f(x)- \eps^{-1}\alpha(x) I_{\mu_\eps}(t)\right)\partial_x|p_\eps|
\\
 + \frac{\sigma^2}{2}\left(\frac{\gamma''(w_\eps)}{\eps}+\frac{\gamma'''(w_\eps)}{\gamma'(w_\eps)}-\frac{\gamma''(w_\eps)^2}{\gamma'(w_\eps)^2}\right) |p_\eps|^3
 \\ 
 + \sigma^2\left(\frac{\gamma'(w_\eps)}{\eps}+\frac{\gamma''(w_\eps)}{\gamma'(w_\eps)}\right) p_\eps \cdot\partial_x |p_\eps|.
\end{multline*}
Let us now specify the map $\gamma: x\mapsto - x^2+2F(T)^2$. Recall that, by definition, we have $F(T)\leq w_\eps$, and, moreover
$$
\frac{\gamma''(w_\eps)}{\eps}+\frac{\gamma'''(w_\eps)}{\gamma'(w_\eps)}-\frac{\gamma''(w_\eps)^2}{\gamma'(w_\eps)^2} \leq -\frac{2}{\eps}-\frac{1}{w_\eps^2}\leq -2\eps^{-1}.
$$
We also recall that $f$ and $\alpha$ have linear growth at infinity, thus all their derivatives are bounded. Altogether, we have that:
\begin{multline*}
\partial_t|p_\eps|-\frac{\sigma^2}{2}\partial^2_{xx}|p_\eps|
\leq \frac{C_M(1+\eps)}{2F(T)}+\frac{C_M(1+\eps)}{2F(T)^2}|p_\eps|
+C_M(1+\eps^{-1})|p_\eps|- \sigma^2\eps^{-1} |p_\eps|^3
\\
-  \left( f(x)- \eps^{-1}\alpha(x) I_{\mu_\eps}(t)\right)\partial_x|p_\eps|
 + \sigma^2\left(\frac{\gamma'(w_\eps)}{\eps}+\frac{\gamma''(w_\eps)}{\gamma'(w_\eps)}\right) p_\eps \cdot\partial_x |p_\eps|.
\end{multline*}
Assuming that $\eps<1$ (consistent with the large connectivity regime we are interested in), there exists a constant $C_M(T)$ such that:
\begin{multline*}
\partial_t|p_\eps|-\frac{\sigma^2}{2}\partial^2_{xx}|p_\eps|
\leq C_M(T)(1+|p_\eps|)
- \frac{\sigma^2}2|p_\eps|^3
+\eps^{-1}\left(C_M|p_\eps|- \frac{\sigma^2}2 |p_\eps|^3\right)
\\
-  \left( f(x)- \eps^{-1}\alpha(x) I_{\mu_\eps}(t)\right)\partial_x|p_\eps|
 + \sigma^2\left(\frac{\gamma'(w_\eps)}{\eps}+\frac{\gamma''(w_\eps)}{\gamma'(w_\eps)}\right) p_\eps \cdot\partial_x |p_\eps|. 
\end{multline*}
Noting that the expression on the first line is a third-degree polynomial in $\vert p_\eps\vert$, with coefficients either of order 1 or $\eps^{-1}$ and depending on $T$, there exists a large enough positive constant $\theta(T)$ such that
\begin{multline*}
\partial_t|p_\eps|-\frac{\sigma^2}{2}\partial^2_{xx}|p_\eps|
\leq -\frac{\sigma^2}{2}\left(|p_\eps|-\theta(T)\right)^3(1+\eps^{-1})
\\
-  \left( f(x)- \eps^{-1}\alpha(x) I_{\mu_\eps}(t)\right)\partial_x|p_\eps|
 + \sigma^2\left(\frac{\gamma'(w_\eps)}{\eps}+\frac{\gamma''(w_\eps)}{\gamma'(w_\eps)}\right) p_\eps \cdot\partial_x |p_\eps|.
\end{multline*}
Noticing that the spatially homogeneous function $y_\eps(t) = \sqrt{\frac{\eps}{\sigma^2(1+\eps)t}}+\theta(T),$ which is such that $y(0)=\infty$, is a solution of the equation:
\begin{multline*}
\partial_ty(t,x)-\frac{\sigma^2}{2}\partial^2_{xx}y(t,x)
= -\frac{\sigma^2}{2}\left(y(t,x)-\theta(T)\right)^3(1+\eps^{-1})
\\
-  \left( f(x)- \eps^{-1}\alpha(x) I_{\mu_\eps}(t)\right)\partial_xy(t,x)
 + \sigma^2\left(\frac{\gamma'(w_\eps)}{\eps}+\frac{\gamma''(w_\eps)}{\gamma'(w_\eps)}\right) y(t,x) \cdot\partial_x y(t,x),
\end{multline*}
we conclude that 
\[
|p_\eps(t,x)|\leq  \sqrt{\frac{\eps}{\sigma^2(1+\eps)t}}+\theta(T),\qquad 0<t\leq T. 
\]
\end{proof}

\end{document}
